% ============================================================================
% 02-preliminaries.tex — DRAFTED IN FULL
% ============================================================================
\section{Preliminaries and notation}\label{sec:prelim}

\subsection{Operators}

Throughout, $z$ is the affine coordinate, $D = \frac{d}{dz}$, and
$\thetaop = zD$ is the Euler operator. We work with linear differential
operators with coefficients in $\Q[z]$ (occasionally $\Q(z)$ when a monic
normalization is needed), acting on formal power series and on multivalued
analytic functions on $\mathbb{P}^1$ minus the singular locus.

An order-2 operator will be written in $\thetaop$-form as
\begin{equation}\label{eq:L2-theta-form}
  \Ltwo \;=\; P_2\,\thetaop^2 + P_1\,\thetaop + P_0,
  \qquad P_i\in\Q[z],
\end{equation}
and an order-3 operator as
$\Lthree = c_3\,\thetaop^3 + c_2\,\thetaop^2 + c_1\,\thetaop + c_0$ with
$c_i \in \Q[z]$. For a monic operator in $D$-form,
$L = D^2 + pD + q$ with $p,q$ rational functions, the \emph{symmetric square}
$\Symtwo(L)$ is the unique monic order-3 operator annihilating all products
$y_iy_j$ of solutions of $L$; explicitly
\begin{equation}\label{eq:sym2-formula}
  \Symtwo(D^2+pD+q) \;=\;
  D^3 + 3p\,D^2 + \bigl(2p^2 + p' + 4q\bigr)D + \bigl(4pq + 2q'\bigr).
\end{equation}
Formula~\eqref{eq:sym2-formula} is validated in the formalization against two
examples derived from first principles (\S\ref{sec:formalization}): the
harmonic oscillator, $\Symtwo(D^2+1) = D^3+4D$ (solution products spanned by
$1, \cos 2t, \sin 2t$), and $\Symtwo(D^2+D) = D^3+3D^2+2D$ (solution products
$1, e^{-t}, e^{-2t}$).

\subsection{Riemann schemes and Fuchsian invariants}

For a Fuchsian operator $L$ of order $n$ on $\mathbb{P}^1$, the
\emph{Riemann scheme} records the local exponents of $L$ at each singular
point. The Fuchs relation states that the sum of all exponents, over all
singular points, equals $\binom{n}{2}(s-2)$ where $s$ is the number of
singular points; for an order-3 operator with four singular points the
required total is $6$. A point of \emph{maximal unipotent monodromy} (MUM) is
one whose exponents all coincide (here: $0,0,0$ at $z=0$), so that the local
monodromy is a single Jordan block. We write $W(L)$ for the Wronskian of a
fundamental system of solutions of $L$.

\subsection{Cooper's template and the three sporadic sequences}

Cooper~\cite{Cooper2012} introduced three sequences $s_7$, $s_{10}$, $s_{18}$
satisfying order-3 Ap\'ery-like recurrences; Gorodetsky~\cite{Gorodetsky2023}
subsumed them in the template~\eqref{eq:cooper-template}, whose
$\thetaop$-form coefficients, after expanding and moving powers of $z$ to the
left (the inner factors have constant coefficients, so no commutator terms
arise), are
\begin{equation}\label{eq:cooper-coeffs}
\begin{aligned}
  c_3 &= 1 - 2a\,z + c\,z^2, &
  c_2 &= -3a\,z + 3c\,z^2, \\
  c_1 &= -(a+2b)\,z + (3c+d)\,z^2, &
  c_0 &= -b\,z + (c+d)\,z^2.
\end{aligned}
\end{equation}
The holomorphic solution of $\Lthree(a,b,c,d)$ at $z=0$ normalized by
$u(0)=1$ has coefficients satisfying Cooper's three-term recurrence; for
$(a,b,c,d)=(13,4,-27,3)$ this is $s_7$ (OEIS A183204: $1, 4, 48, 760,
\dots$), for $(6,2,-64,4)$ it is $s_{10}$ (OEIS A005260, the sums
$\sum_k\binom{n}{k}^4$), and for $(14,6,192,-12)$ it is $s_{18}$
\cite{Cooper2012,Gorodetsky2023}. The parameter values are transcribed from
the literature and are pinned, with source citations attached to the
definitions, in the formal development (\S\ref{sec:formalization}).

\subsection{Lattices}

A \emph{lattice} is a free $\Z$-module of finite rank with a nondegenerate
symmetric bilinear form; it is \emph{even} if $\lat{x,x}\in 2\Z$ for all $x$.
$\U$ denotes the hyperbolic plane
$\left[\begin{smallmatrix}0&1\\1&0\end{smallmatrix}\right]$ and $\lat{m}$ the
rank-1 lattice with Gram matrix $(m)$. The \emph{discriminant group} of a
lattice $T$ is $T^\vee/T$, carrying its discriminant quadratic form
$q\colon T^\vee/T\to\Q/2\Z$. Two lattices are isometric over $\Z$ if their
Gram matrices are conjugate under $\GL_n(\Z)$.

For a K3 surface $X$, $\NS(X)\subset H^2(X,\Z)$ is the N\'eron--Severi
lattice, of rank $\rho(X)$, and the \emph{transcendental lattice} $T(X)$ is
its orthogonal complement, of rank $22-\rho(X)$. Following
Dolgachev~\cite{Dolgachev1996}, for $M$ a fixed even lattice of signature
$(1,t)$, an \emph{$M$-polarized} K3 surface is one with a primitive lattice
embedding $M\hookrightarrow \NS(X)$; for the rank-19 lattices
$M_n = \U\oplus E_8^2\oplus\lat{-2n}$, Dolgachev computes
$(M_n)^{\perp} = \U\oplus\lat{2n}$ (\cite{Dolgachev1996}, \S7) and proves
that the coarse moduli space of $M_n$-polarized K3 surfaces is the Fricke
modular curve $H/\Gamma_0(n)+$ (\cite{Dolgachev1996}, Thm 7.1).

% ----------------------------------------------------------------------------
% Generated-by: Fable 5 (Stream 1 paper agent) | Verified-by: eq (coeffs) vs
% Agora/Sequences/PartnerOperators.lean cooperC3..C0; params vs
% CooperRecurrences.lean; Dolgachev statements vs Stream 2 sources-read brief
% | Reviewed-by: pending T0 (Xavier)
% ----------------------------------------------------------------------------
